\chapter{IMPLEMENTATION DETAILS}
\section{Methodology Implementation}
\subsection{Geometric camera Callibration}
Camera calibration has been completed for all three cameras, using [N] checkerboard image pairs per camera. Table~\ref{tab:reprojection_error} reports the resulting mean reprojection error per camera, all within acceptable bounds for accurate triangulation.
\begin{table}[ht]
\centering
\caption{Camera Reprojection Errors}
\label{tab:reprojection_error}
\begin{tabular}{lc}
\toprule
\textbf{Camera} & \textbf{Projection Error (px)} \\
\midrule
Camera 1 (Front) & X.XX \\
Camera 2 (Left)  & X.XX \\
Camera 3 (Right) & X.XX \\
\bottomrule
\end{tabular}
\end{table}

\subsection{Parallel Multi-Modal Keypoint Extraction}
 MediaPipe Pose Landmarker and MediaPipe Hands have been integrated and tested across all three synchronized camera streams. Live landmark overlays confirming correct detection across all views are shown in Figure 4-6(same ss). The pipeline retains only the shoulder (landmarks 11, 12), elbow (13, 14), and wrist (15, 16) points from MediaPipe Pose's 33-point topology, alongside all 21 points from MediaPipe Hands per detected hand.
 % need to insert an screenshot of a real person's arm/hand, captured live by your webcam(s), with MediaPipe's landmark dots and skeleton connections drawn directly on top of the video frame — proving detection is actually running on your real setup, not a mockup. we also have demo video so maybe not needed.idkkk
\subsubsection{Confidence-Weighted Triangulation}
Confidence-weighted triangulation (Equation~\ref{eq:4-3}) has been implemented and validated across the full three-camera configuration. Reconstruction accuracy was assessed against a physically measured reference segment (e.g., [XX] cm forearm distance), yielding a mean absolute error of [X.XX] cm across [N] trials.

\subsection{Vector Proportional Retargeting}
Vector-proportional retargeting has been implemented, normalizing captured human joint vectors to the Kuka IIWA URDF model's link lengths as described in Equation~\ref{eq:4-4}, preventing joint hyper-extension during simulation.

\subsection{Articulated Kinematic Setup within PyBullet}
The PyBullet simulation environment has been implemented using the Kuka IIWA URDF manipulator model, rendering physically accurate joint movement in real time. Figure 4-7(same ss) shows the simulated arm replicating captured human motion.
% need to insert screenshot of the actual PyBullet simulation window, open and rendering, with the Kuka IIWA robotic arm posed in a way that visibly matches (or is clearly derived from) a captured human arm position — proving the IK mapping from human motion to robot joints is actually functioning.but again demo vid so idkkkk
\subsection{Damped Least Squares Inverse Kinematics}
The DLS-based IK solver(Equation~\ref{eq:4-5}) has been implemented and validated, successfully mapping reconstructed $3$D wrist positions to joint configurations without encountering singularity-related failures during testing. 

\subsection{Automated Data Factory Serialization}
Automated CSV dataset serialization has been implemented, producing per-frame records of 3D joint coordinates, computed joint angles, and timestamps.

\subsection{Sequential Noise Mitigation Integration}

A real-time analytical smoothing stage was integrated into the live capture pipeline using the One Euro Filter. The filter was applied immediately after triangulation and prior to inverse kinematics computation. This reduced high-frequency positional jitter without affecting frame rate performance. The filtered 3D coordinates were simultaneously forwarded to the IK solver for simulation and recorded in the dataset serialization module.

Following data acquisition, an offline refinement stage was implemented using a lightweight stacked LSTM model. The network processes short temporal windows of recorded joint coordinates to further enhance motion continuity and reduce residual micro-oscillations. This post-processing step operates independently from the real-time pipeline and produces a curated dataset suitable for robotic imitation learning and further experimental evaluation.